% !TEX root = ../discretemath.tex

\section{Пример для асимметричной версии локальной леммы Ловаса}

\subsection{Напоминание: асимметричная форма}

\begin{Theorem}{Асимметричная форма ЛЛЛ (формулировка)}{}
	Пусть $\Gamma$ — граф зависимостей семейства $\{A_i\}$, и существует $f:\{A_i\}\to[0,1)$ такая, что для каждого $i$
	\[
	\Pr(A_i)\le f(A_i)\prod_{k\in\Gamma^+(i)}\bigl(1-f(A_k)\bigr).
	\]
	Тогда $\Pr\!\left(\bigwedge_i\overline{A_i}\right)\ge\prod_i\bigl(1-f(A_i)\bigr)>0$. (Доказательство — в билете о локальной лемме Ловаса.)
\end{Theorem}

\subsection{Нижняя оценка на $R(3,k)$}

\begin{Example}{Нижняя оценка на $R(3,k)$}{}
	Рассмотрим случайный граф $G$ на $N$ вершинах: каждое ребро проводим независимо с вероятностью $q$. Цель — показать, что при подходящих $N,q$ с положительной вероятностью в $G$ нет ни треугольника, ни независимого множества размера $k$, откуда $R(3,k)>N$.
\end{Example}

Введём два типа «плохих» событий:
\[
A_S=\{\text{на }S\text{ полный граф }K_3\},\ |S|=3;
\qquad
B_T=\{\text{на }T\text{ пустой граф}\},\ |T|=k.
\]
Тогда
\[
\Pr(A_S)=q^3
\quad(\text{три ребра треугольника}),
\qquad
\Pr(B_T)=(1-q)^{\binom{k}{2}}
\quad(\text{все }\tbinom{k}{2}\text{ рёбер отсутствуют}).
\]

\begin{figure}[H]
	\centering
	\begin{tikzpicture}[font=\small,
		v/.style={circle, draw, fill=black!80, inner sep=0pt, minimum size=4pt}]
		% A_S: треугольник
		\begin{scope}
			\node[v] (a) at (0,0) {}; \node[v] (b) at (1.1,0) {}; \node[v] (c) at (0.55,0.9) {};
			\draw[thick] (a)--(b)--(c)--(a);
			\node at (0.55,-0.55) {$A_S$: треугольник ($q^3$)};
		\end{scope}
		% B_T: пустое множество
		\begin{scope}[xshift=4cm]
			\foreach \ang in {90,150,210,270,330,30}{ \node[v] at (\ang:0.7) {}; }
			\node at (0,-1.2) {$B_T$: нет рёбер ($(1-q)^{\binom{k}{2}}$)};
		\end{scope}
	\end{tikzpicture}
	\caption{Два типа запрещённых конфигураций. Их вероятности различаются на много порядков, поэтому единый параметр симметричной леммы не подходит — применяется асимметричная форма с разными $f(A_S)=x$ и $f(B_T)=y$.}
\end{figure}

\subsubsection*{Граф зависимостей}

События зависимы, лишь когда соответствующие множества вершин имеют общую \emph{пару} вершин (общее потенциальное ребро). Отсюда оценки числа зависимостей:
\[
A_S \text{ зависит от } \le 3(N-3) \text{ событий } A_{S'},
\qquad
A_S \text{ зависит от } \le 3\binom{N-2}{k-2} \text{ событий } B_T,
\]
\[
B_T \text{ зависит от } \le \binom{k}{2}(N-2) \text{ событий } A_S,
\qquad
B_T \text{ зависит от } \le \binom{k}{2}\binom{N-2}{k-2} \text{ событий } B_{T'}.
\]

\subsubsection*{Применение асимметричной ЛЛЛ}

Положим $f(A_S)=x$, $f(B_T)=y$. Достаточно подобрать $x,y\in[0,1)$ так, чтобы
\[
q^3\le x(1-x)^{3(N-3)}(1-y)^{3\binom{N-2}{k-2}},
\]
\[
(1-q)^{\binom{k}{2}}\le y(1-x)^{\binom{k}{2}(N-2)}(1-y)^{\binom{k}{2}\binom{N-2}{k-2}}.
\]

Если такие параметры существуют, то с положительной вероятностью в $G$ нет ни треугольника, ни независимого множества размера $k$, и потому
\[
R(3,k)>N.
\]

\begin{Remark}{}{}
	При оптимальном выборе $q\sim c/\sqrt{N}$ и согласованных $x,y$ метод даёт
	\[
	R(3,k)\ge\Omega\!\left(\left(\frac{k}{\log k}\right)^2\right).
	\]
	Это существенно сильнее, чем даёт прямой вероятностный метод без локальной леммы, и иллюстрирует, зачем нужна именно асимметричная форма: вклад редких, но «крупных» событий $B_T$ балансируется отдельной функцией $y$.
\end{Remark}
